% Chapter 13: Monte Carlo Greeks.

A price is one number; a hedge needs its derivatives. The delta and the gamma decide the
hedge and how often to rebalance it (\cref{ch:18-hedging}), the vega and the rhos measure
the exposure to volatility and rates, and a risk system recomputes all of them every day
for every position. When the price comes from a simulation, so must the Greeks. The course
compares three ways to get them: \emph{bumping} an input and pricing again, differentiating
the simulated payoff path by path (the course's \emph{analytical differencing}, called the
\emph{pathwise} method in the literature), and differentiating the probability density of
the simulated state (the \emph{likelihood-ratio} method). The gamma needs mixtures of the
last two, and books with many inputs need the \emph{adjoint} method, which computes every
first-order Greek of a path at a small multiple of the cost of its price.

The three approaches trade bias, variance and regularity in different ways. Bumping is
biased, and done naively its variance explodes as the bump shrinks. The pathwise method is
unbiased with a small variance, but it needs a payoff that is Lipschitz in the parameter.
The likelihood-ratio method needs only a smooth density, but its variance is usually larger,
and it grows with the number of time steps and as $\sigma\sqrt T\to0$. The code is
\coderef{code/src/mfe230d/greeks.py}.

Throughout, a Greek is the derivative $\alpha'(\theta)$ of $\alpha(\theta)=\E[Y(\theta)]$,
where $Y(\theta)$ is a discounted payoff and $\theta$ one input. The running example is the
test case $S=K=100$, $T=1$, $r=1\%$, $y=2\%$, $\sigma=20\%$, with
\[
  S_T=S\,G,\qquad G=\exp\big(\nu T+\sigma\sqrt T\,Z\big),\qquad \nu=r-y-\tfrac{\sigma^2}2,\qquad Z\sim\mathcal N(0,1)
\]
under $\Q$, call price $7.36429$, and Black--Scholes Greeks (\texttt{bs.greeks},
\cref{prop:bsf-greeks}) $\Delta=e^{-yT}N(d_+)=0.509643$, $\Gamma=0.0195277$, vega
$\mathcal V=39.0554$, theta $\partial_tV=-3.32226$ per year, $\rho_r=43.6000$ and
$\rho_y=-50.9643$. The cash-or-nothing digital call pays $\1_{\{S_T>K\}}$; its price is
$e^{-rT}N(d_-)=0.436000$ and its delta $e^{-rT}N'(d_-)/(S\sigma\sqrt T)=0.0195277$, equal to
the gamma of the call because $S\,\Gamma_{\mathrm{call}}=K\,\Delta_{\mathrm{dig}}$ (from
$Se^{-yT}N'(d_+)=Ke^{-rT}N'(d_-)$).

\section{Finite differences}\label{sec:gr-fd}

Bumping is the natural first approach: it treats the pricer as a black box and works for any
model. The course bumps the spot by a relative amount, $S_\pm=S(1\pm\varepsilon)$, and the
other inputs by absolute amounts (one basis point for $r$ and $y$, one volatility point for
$\sigma$). With $C_0$ and $C_\pm$ the Monte Carlo prices at $S$ and $S_\pm$, the forward,
backward and central deltas and the central gamma are
\begin{equation}\label{eq:gr-course-fd}
  \Delta_+=\frac{C_+-C_0}{\varepsilon S},\qquad \Delta_-=\frac{C_0-C_-}{\varepsilon S},\qquad
  \Delta_c=\frac{C_+-C_-}{2\varepsilon S},\qquad \Gamma_c=\frac{C_++C_--2C_0}{(\varepsilon S)^2}.
\end{equation}
Two choices matter: the size $h=\varepsilon S$ of the bump, and whether the three prices use
the same random numbers (\emph{common random numbers}, CRN; \cref{prop:cv-crn}) or
independent ones. The next proposition, for a generic input $\theta$, quantifies both.

\begin{proposition}[Finite differences: bias, variance and the optimal bump]\label{prop:gr-finite-difference}
Let $I$ be an open interval and, for $\theta\in I$, let $Y(\theta)=g(\theta,U)$ with $U$ a
random vector and $\E[Y(\theta)^2]<\infty$; write $\alpha(\theta)=\E[Y(\theta)]$ and
$\sigma^2(\theta)=\Var(Y(\theta))$. Fix $\theta$ and $h>0$ with $[\theta-h,\theta+h]\subset I$,
and let $\hat\Delta^F_{n}$, $\hat\Delta^C_{n}$ and $\hat\Gamma_{n}$ be the averages over $n$
independent replications of the forward difference $(Y(\theta+h)-Y(\theta))/h$, the central
difference $(Y(\theta+h)-Y(\theta-h))/(2h)$ and the second difference
$(Y(\theta+h)-2Y(\theta)+Y(\theta-h))/h^2$.
\begin{enumerate}[label=(\roman*)]
  \item (\emph{Bias}, whatever the sampling.) If $\alpha$ is $C^2$ near $\theta$, then
    $\E\hat\Delta^F_n-\alpha'(\theta)=\frac12\alpha''(\theta)h+o(h)$. If $\alpha$ is $C^3$, then
    $\E\hat\Delta^C_n-\alpha'(\theta)=\frac16\alpha'''(\theta)h^2+o(h^2)$. If $\alpha$ is $C^4$, then
    $\E\hat\Gamma_n-\alpha''(\theta)=\frac1{12}\alpha^{(4)}(\theta)h^2+o(h^2)$.
  \item (\emph{Independent sampling}: each evaluation within a replication uses its own
    independent copy of $U$.) Then
    \begin{gather*}
      \Var\hat\Delta^F_n=\frac{\sigma^2(\theta+h)+\sigma^2(\theta)}{nh^2},\qquad
      \Var\hat\Delta^C_n=\frac{\sigma^2(\theta+h)+\sigma^2(\theta-h)}{4nh^2},\\
      \Var\hat\Gamma_n=\frac{\sigma^2(\theta+h)+4\sigma^2(\theta)+\sigma^2(\theta-h)}{nh^4}.
    \end{gather*}
    If $\sigma^2$ is continuous and positive at $\theta$, these are of exact order
    $1/(nh^2)$, $1/(nh^2)$ and $1/(nh^4)$.
  \item (\emph{Common random numbers}: one $U$ per replication.) Then
    $\Var\hat\Delta^C_n=\Var\big(Y(\theta+h)-Y(\theta-h)\big)/(4nh^2)$, and similarly for the
    others.
    \begin{enumerate}[label=(\alph*)]
      \item If there are $\delta>0$ and a random variable $\kappa$ with $\E\kappa^2<\infty$
        and $|Y(\theta_1)-Y(\theta_2)|\le\kappa|\theta_1-\theta_2|$ almost surely for all
        $\theta_1,\theta_2\in(\theta-\delta,\theta+\delta)$ (a \emph{Lipschitz} payoff),
        then $\Var\hat\Delta^F_n$ and $\Var\hat\Delta^C_n$ are at most $\E\kappa^2/n$ for every
        $h<\delta$.
      \item If $\E\big[(Y(\theta_1)-Y(\theta_2))^2\big]\le c|\theta_1-\theta_2|$ for
        $\theta_1,\theta_2$ near $\theta$ (a payoff with a jump, \cref{exa:gr-digital}), then
        $\Var\hat\Delta^F_n\le c/(nh)$ and $\Var\hat\Delta^C_n\le c/(2nh)$.
      \item If $|Y(\theta+h)-2Y(\theta)+Y(\theta-h)|\le h\kappa\1_{A_h}$ for a random $\kappa\ge0$
        and events $A_h$ with $\E[\kappa^2\1_{A_h}]\le ch$ (a payoff with a kink,
        \cref{rem:gr-kernel}), then $\Var\hat\Gamma_n\le c/(nh)$.
    \end{enumerate}
  \item (\emph{Optimal bump}.) Suppose that, as $h\downarrow0$, the bias is
    $bh^\beta(1+o(1))$ and $n$ times the variance is $vh^{-\gamma}(1+o(1))$, with constants
    $b\ne0$, $v>0$, $\beta>0$, $\gamma>0$ that do not depend on $n$. Let
    $c^*=\big(\gamma v/(2\beta b^2)\big)^{1/(2\beta+\gamma)}$. For every sequence
    $h_n\to0$,
    \[
      \liminf_{n\to\infty}\,n^{2\beta/(2\beta+\gamma)}\,\E\big[(\hat\Delta_n-\alpha'(\theta))^2\big]
      \ \ge\ \Big(1+\frac{2\beta}{\gamma}\Big)b^2(c^*)^{2\beta},
    \]
    with equality (and a limit) for $h_n=c^*n^{-1/(2\beta+\gamma)}$; here $\hat\Delta_n$ is
    the estimator built with the bump $h_n$, and $\alpha'(\theta)$ is replaced by
    $\alpha''(\theta)$ for $\hat\Gamma_n$. In case (iii)(a), the mean squared error is at
    most $b^2h^{2\beta}(1+o(1))+\E\kappa^2/n$, which is $O(1/n)$ for every $h_n=O(n^{-1/(2\beta)})$.
\end{enumerate}
\end{proposition}

\begin{proof}
(i) By Taylor's theorem with the Lagrange remainder,
$\alpha(\theta+h)=\alpha(\theta)+\alpha'(\theta)h+\frac12\alpha''(\xi)h^2$ with $\xi$ between
$\theta$ and $\theta+h$, and $\alpha''(\xi)\to\alpha''(\theta)$ by continuity. For the central
difference, expand $\alpha(\theta\pm h)$ to third order: the even terms cancel and
$\alpha(\theta+h)-\alpha(\theta-h)=2\alpha'(\theta)h+\frac16(\alpha'''(\xi_+)+\alpha'''(\xi_-))h^3$.
For the second difference, expand to fourth order: the odd terms cancel and
$\alpha(\theta+h)+\alpha(\theta-h)-2\alpha(\theta)=\alpha''(\theta)h^2+\frac1{24}(\alpha^{(4)}(\xi_+)+\alpha^{(4)}(\xi_-))h^4$.
In every case the expectation of the estimator is the same difference of $\alpha$, by linearity,
whatever the dependence between the evaluations.

(ii) The replications are independent and, within one, the evaluations are independent, so
the variances add: $\Var(aY_1+bY_2+cY_3)=a^2\sigma_1^2+b^2\sigma_2^2+c^2\sigma_3^2$, and averaging
over $n$ replications divides by $n$.

(iii) The first claim is the variance of an average of i.i.d.\ terms. (a) With
$|Y(\theta+h)-Y(\theta-h)|\le2h\kappa$, the variance of one central quotient is at most its
second moment, $\le\E\kappa^2$; the forward quotient is the same. (b) Likewise, the variance
of one quotient is at most $ch/h^2$ (forward) or $2ch/(4h^2)$ (central). (c) The second moment
of one second difference is at most $h^2\E[\kappa^2\1_{A_h}]/h^4\le c/h$.

(iv) The mean squared error is the squared bias plus the variance
(\cref{prop:mc-mse}(i)). Write $h_n=c_nn^{-1/(2\beta+\gamma)}$. Then
$n^{2\beta/(2\beta+\gamma)}h_n^{2\beta}=c_n^{2\beta}$ and
$n^{2\beta/(2\beta+\gamma)}h_n^{-\gamma}/n=c_n^{-\gamma}$, so
\[
  n^{2\beta/(2\beta+\gamma)}\,\mathrm{MSE}=b^2c_n^{2\beta}(1+o(1))+vc_n^{-\gamma}(1+o(1))
  \ \ge\ (1-\eta)\,\phi(c_n)\ \ge\ (1-\eta)\,\phi(c^*)
\]
for every $\eta\in(0,1)$ and $n$ large, where $\phi(c)=b^2c^{2\beta}+vc^{-\gamma}$. As in the
proof of \cref{prop:mc-mse}(iii), $\phi$ decreases on $(0,c^*]$ and increases on $[c^*,\infty)$, so $c^*$
is its global minimiser, and $vc^{*-\gamma}=\frac{2\beta}{\gamma}b^2c^{*2\beta}$ gives
$\phi(c^*)=(1+2\beta/\gamma)b^2c^{*2\beta}$. With $c_n=c^*$ the two $o(1)$ terms vanish in the
limit. The last claim combines the bias with the bound of (iii)(a).
\end{proof}

The proposition gives the familiar table of rates
\citep[Sections~7.1.1--7.1.2]{glasserman-mc}, due to \textcite{glynn-1989} and
\textcite{lecuyer-perron}:
\begin{center}
\small
\begin{tabular}{@{}lccccc@{}}
  \toprule
  & $\beta$ & $\gamma$ & best bump & best MSE & RMSE \\
  \midrule
  forward, independent & $1$ & $2$ & $n^{-1/4}$ & $n^{-1/2}$ & $n^{-1/4}$ \\
  central, independent & $2$ & $2$ & $n^{-1/6}$ & $n^{-2/3}$ & $n^{-1/3}$ \\
  forward, CRN, payoff with a jump & $1$ & $1$ & $n^{-1/3}$ & $n^{-2/3}$ & $n^{-1/3}$ \\
  central, CRN, payoff with a jump & $2$ & $1$ & $n^{-1/5}$ & $n^{-4/5}$ & $n^{-2/5}$ \\
  central or forward, CRN, Lipschitz payoff & $1$ or $2$ & --- & small enough & $n^{-1}$ & $n^{-1/2}$ \\
  \midrule
  second difference, independent & $2$ & $4$ & $n^{-1/8}$ & $n^{-1/2}$ & $n^{-1/4}$ \\
  second difference, CRN, payoff with a kink & $2$ & $1$ & $n^{-1/5}$ & $n^{-4/5}$ & $n^{-2/5}$ \\
  \bottomrule
\end{tabular}
\end{center}
In the table, only common random numbers with a Lipschitz payoff reach the Monte Carlo rate
$n^{-1/2}$; higher-order differences come closer without reaching it (\exref{ex:gr-extrapolation}).
\textcite{lecuyer-perron} prove the rate $n^{-1/2}$ under the conditions that also make the
pathwise method of \cref{sec:gr-pathwise} valid. The last line is
\cref{rem:gr-kernel}. Everything assumes $b\ne0$. If $b=0$ the bias is of higher order and
the rates improve; if $b$ is merely small, the asymptotic rates set in only for very large $n$
(\cref{fig:gr-bump}(c)).

\begin{example}[The delta of a digital]\label{exa:gr-digital}
Let $Y(S)=e^{-rT}\1_{\{SG>K\}}$ with common draws. Then
$Y(S+h)-Y(S-h)=e^{-rT}\1_{\{K/(S+h)<G\le K/(S-h)\}}$, so
\[
  \E\big[(Y(S+h)-Y(S-h))^2\big]=e^{-2rT}\big(N(z_{-h})-N(z_{+h})\big),\qquad
  z_{\pm h}=\frac{\log(K/(S\pm h))-\nu T}{\sigma\sqrt T}.
\]
By the mean value theorem,
\[
  N(z_{-h})-N(z_{+h})=N'(\xi_h)(z_{-h}-z_{+h}),\qquad
  z_{-h}-z_{+h}=\frac{1}{\sigma\sqrt T}\log\frac{S+h}{S-h}=\frac{2h}{S\sigma\sqrt T}\big(1+o(1)\big),
\]
with $\xi_h\to-d_-$. So the second moment is $2he^{-rT}\Delta_{\mathrm{dig}}(1+o(1))$: case (iii)(b)
with $\gamma=1$ and $v=e^{-rT}\Delta_{\mathrm{dig}}/2$. The bias of the central difference is
$\frac16\alpha'''(S)h^2$, with $\alpha'''(S)=-4.82\times10^{-5}$ for the test case, and (iv)
gives $h^*=2.06$ at $n=10^6$, a $2\%$ bump; minimising the exact mean squared error (bias from
the closed form, variance by quadrature, \texttt{greeks.optimal\_bump}) gives $h^*=2.07$ and a
root mean squared error of $7.39\times10^{-5}$, $0.38\%$ of the delta. Without common random
numbers the best is $1.35\times10^{-4}$, and the likelihood-ratio estimator of
\cref{sec:gr-lr} has standard error $0.0290/\sqrt n=2.90\times10^{-5}$ with no bias at all: at
equal accuracy it needs $6.5$ times fewer draws than the best bump.
\end{example}

\begin{remark}[The bumped gamma is a kernel estimator]\label{rem:gr-kernel}
For the call with common draws, $(x+h)^+-2x^++(x-h)^+=(h-|x|)^+$ gives, sample by sample,
\[
  \frac{Y(S+h)-2Y(S)+Y(S-h)}{h^2}=e^{-rT}G\,k_h\big(S-K/G\big),\qquad k_h(u)=\frac{(h-|u|)^+}{h^2},
\]
and $\Delta_c$ is $e^{-rT}G$ times the ramp $\int_{-\infty}^{S-K/G}\1_{\{|u|<h\}}\,\dd u/(2h)$. Each draw has a
\emph{critical spot} $K/G$, the initial price at which it finishes at the money. The bumped
gamma $\hat\Gamma_n=\frac1n\sum_ie^{-rT}G_ik_h(S-K/G_i)$ is a kernel estimator, with the
triangular kernel $k_h$ of width $2h$ and integral $1$, of the density of the critical spots
weighted by $e^{-rT}G$, which is the gamma. Since the tent is at most $h$, each second
difference is at most $he^{-rT}G\1_{A_h}$ with $A_h=\{|S-K/G|<h\}$, and
$\E[e^{-2rT}G^2\1_{A_h}]=O(h)$ because $K/G$ has a bounded density near $S$ and $G\le2K/S$ on
$A_h$ for $h<S/2$: \cref{prop:gr-finite-difference}(iii)(c) bounds the variance by a multiple of
$1/(nh)$, and the bias is $O(h^2)$, the last line of the table. The course's region-by-region analysis is this
decomposition: regions II and III, where $K/(1+\varepsilon)<S_T<K/(1-\varepsilon)$, are the two
sides of the tent, and outside them a draw contributes nothing to $\Gamma_c$. The number of
draws in the tent is binomial, with mean $n\,\Q(|S-K/G|<h)$: $4.0$ for the course's $n=1024$
and $\varepsilon=0.001$, $40$ for $\varepsilon=0.01$ and $388$ for $\varepsilon=0.1$. With four
draws the estimate is a handful of tents and zero in between, the ``butterflies'' of
\cref{fig:gr-course}(b). The deltas are sums of ramps, which become steps as $h\to0$: the
``$1/n$ quantization'' of the course.
\end{remark}

\begin{figure}[ht]
  \centering
  \includegraphics[width=\textwidth]{ch13-bump}
  \caption{Finite differences for the test case, exact curves (bias from the closed forms,
    variance by quadrature). (a) Root mean squared error of the call delta against the relative
    bump, $n=10^4$; dots: $200$ independent runs each (our implementation, \texttt{default\_rng},
    seeds $100$--$103$). With CRN the error is flat at the pathwise standard error $0.0058$.
    (b) Error relative to the exact value for the central digital delta and the second-difference
    call gamma, $n=10^4$; with CRN the two curves almost coincide (their exact values are equal
    and both variances are of order $1/(nh)$). (c) Error at the optimal bump against $n$ for the
    digital delta; reference slopes $-1/4$, $-1/3$, $-2/5$. The forward curves approach their
    slopes slowly because the digital's $\alpha''=-4.9\times10^{-5}$ is small at the test case.}
  \label{fig:gr-bump}
\end{figure}

\Cref{fig:gr-bump}(a) shows the two regimes for the call delta. With independent draws the
error grows like $1/\varepsilon$ for small bumps, and the best forward difference ($\varepsilon=4.4\%$)
has error $0.061$. With common random numbers the central difference has error $0.0058$ for
every bump up to $3\%$, ten times smaller, which is worth a hundred times more paths; the
forward difference stays there up to $0.1\%$, after which its bias of order $h$ takes over. A
bump of $10^{-4}$ is as good as a bump of $10^{-2}$: with a Lipschitz payoff and common draws,
small bumps cost nothing. Rounding
does not interfere at these sizes, since in double precision a difference quotient loses about
$10^{-16}|Y|/h$, far below the statistical error. \Cref{fig:gr-course} reproduces the course's
experiment: $n=1024$ draws, reused for every $S_0$. The forward and backward deltas bracket the
central one, $\Delta_+-\Delta_-=\varepsilon S\,\Gamma_c\ge0$ sample by sample because the payoff is
convex, and the gammas are biased for $\varepsilon=0.1$ and useless for $\varepsilon=0.001$.

\begin{figure}[ht]
  \centering
  \includegraphics[width=0.94\textwidth]{ch13-course}
  \caption{The course's shoot-out on the test-case call, $n=1024$ draws reused for every $S_0$
    (our implementation, \texttt{default\_rng}, seed $1024$). (a) Bumped deltas
    \eqref{eq:gr-course-fd}, $\varepsilon=0.01$. (b) Bumped gammas $\Gamma_c$ for
    $\varepsilon=0.1$, $0.01$ and $0.001$ (dots; values above $0.075$ are cut).
    (c) Pathwise and likelihood-ratio deltas, \eqref{eq:gr-pw-bs} and \eqref{eq:gr-lr-bs}.
    (d) The PW-LR and LR-LR gammas of \cref{prop:gr-hybrid-gamma}. Black: Black--Scholes.}
  \label{fig:gr-course}
\end{figure}

\demo{13-mc-greeks.ipynb}{move the bump, the number of draws and the payoff, and watch the
two regimes of the error.}

\section{Pathwise derivatives}\label{sec:gr-pathwise}

Let the bump go to zero with common random numbers. If $\theta\mapsto Y(\theta)$ is
differentiable along each path, the difference quotients converge to the derivative
$Y'(\theta)$ of the payoff along the path, and the question is whether
$\E[Y'(\theta)]=\alpha'(\theta)$: whether the derivative and the expectation can be
interchanged. Almost sure differentiability is not enough (\cref{rem:gr-digital}); what is
needed is a dominating Lipschitz bound.

\begin{theorem}[Pathwise derivatives]\label{thm:gr-pathwise}
Let $I$ be an open interval, $\theta_0\in I$, and $(Y(\theta))_{\theta\in I}$ random variables
with $\E|Y(\theta_0)|<\infty$. Assume
\begin{enumerate}[label=(PW\arabic*)]
  \item for almost every $\omega$, the map $\theta\mapsto Y(\theta,\omega)$ is differentiable
    at $\theta_0$, with derivative $Y'(\theta_0)$;
  \item there are $\delta>0$ and a random variable $\kappa$ with $\E\kappa<\infty$ such that,
    almost surely, $|Y(\theta_1)-Y(\theta_2)|\le\kappa|\theta_1-\theta_2|$ for all
    $\theta_1,\theta_2\in(\theta_0-\delta,\theta_0+\delta)$.
\end{enumerate}
Then $Y(\theta)$ is integrable for $|\theta-\theta_0|<\delta$, $\alpha$ is differentiable at
$\theta_0$, and $\alpha'(\theta_0)=\E[Y'(\theta_0)]$. If moreover $\E\kappa^2<\infty$, then
$\Var Y'(\theta_0)\le\E\kappa^2$ and the common-random-number difference quotients converge to
$Y'(\theta_0)$ in $L^2$ as $h\to0$.

\noindent\emph{Sufficient conditions} \citep[Section~7.2.2, conditions (A1)--(A4)]{glasserman-mc}.
Let $Y(\theta)=f(X(\theta))$ with $X(\theta)$ a random vector in $\R^m$ and $f:\R^m\to\R$
Borel. Then (PW1) and (PW2) hold if: (a) $X(\theta)$ is almost surely differentiable at
$\theta_0$; (b) $\Prob\big(X(\theta_0)\in D_f\big)=1$, where $D_f$ is the set of points at which
$f$ is differentiable; (c) $f$ is Lipschitz, $|f(x)-f(x')|\le L_f\sum_i|x_i-x'_i|$; (d) there
are random variables $\kappa_i$ with $\E\kappa_i<\infty$ and
$|X_i(\theta_1)-X_i(\theta_2)|\le\kappa_i|\theta_1-\theta_2|$ on a neighbourhood of $\theta_0$.
Then $\kappa=L_f\sum_i\kappa_i$ and $Y'(\theta_0)=\nabla f(X(\theta_0))\cdot X'(\theta_0)$. In
particular (PW2) holds, by the mean value theorem, whenever $\theta\mapsto Y(\theta)$ is almost
surely continuously differentiable on the neighbourhood with
$\E\sup_{|\theta-\theta_0|<\delta}|Y'(\theta)|<\infty$.
\end{theorem}

\begin{proof}
For $|\theta-\theta_0|<\delta$, $|Y(\theta)|\le|Y(\theta_0)|+\kappa\delta$ is integrable. Let
$h_k\to0$ with $0<|h_k|<\delta$. By (PW1) the quotients
$Q_k=(Y(\theta_0+h_k)-Y(\theta_0))/h_k$ converge to $Y'(\theta_0)$ almost surely, and by (PW2)
$|Q_k|\le\kappa$, integrable. By dominated convergence,
$(\alpha(\theta_0+h_k)-\alpha(\theta_0))/h_k=\E Q_k\to\E Y'(\theta_0)$; the sequence being
arbitrary, $\alpha'(\theta_0)=\E Y'(\theta_0)$. Also $|Y'(\theta_0)|=\lim|Q_k|\le\kappa$ almost
surely, so $\Var Y'(\theta_0)\le\E\kappa^2$; and $(Q_k-Y'(\theta_0))^2\le4\kappa^2$, so
$\E(Q_k-Y'(\theta_0))^2\to0$ by dominated convergence again; the central quotients are the
average of a forward and a backward one.

For the sufficient conditions: on the almost sure event where (a) and (b) hold, $f$ is
differentiable at $X(\theta_0)$ and $X$ at $\theta_0$, and the chain rule gives (PW1) with the
stated derivative. On the event where (d) holds,
$|f(X(\theta_1))-f(X(\theta_2))|\le L_f\sum_i|X_i(\theta_1)-X_i(\theta_2)|\le L_f\sum_i\kappa_i|\theta_1-\theta_2|$.
Finally, if $\theta\mapsto Y(\theta)$ is $C^1$ on the interval, the mean value theorem gives
$|Y(\theta_1)-Y(\theta_2)|\le\sup|Y'|\,|\theta_1-\theta_2|$.
\end{proof}

Common random numbers and the pathwise method are two views of one thing: under (PW2) with
$\E\kappa^2<\infty$, the bumped quotients with common draws converge in $L^2$ to the pathwise
estimator as $h\to0$, and differ from it in mean by the bias of order $h$ or $h^2$. This is why
the error in \cref{fig:gr-bump}(a) is flat.

\paragraph{Black--Scholes.} Differentiating $S_T=S\exp(\nu T+\sigma\sqrt TZ)$ gives
$\partial_SS_T=G$, $\partial_\sigma S_T=S_T(\sqrt TZ-\sigma T)$, $\partial_rS_T=TS_T$,
$\partial_yS_T=-TS_T$ and $\partial_TS_T=S_T(\nu+\sigma Z/(2\sqrt T))$. For the call,
$f'(s)=\1_{\{s>K\}}$ off the strike, and the pathwise estimators (the course's analytical
differencing) are
\begin{equation}\label{eq:gr-pw-bs}
\begin{aligned}
  \Delta^{\mathrm{pw}}&=e^{-rT}\1_{\{S_T>K\}}G, &\quad
  \mathcal V^{\mathrm{pw}}&=e^{-rT}\1_{\{S_T>K\}}S_T\big(\sqrt TZ-\sigma T\big),\\
  \rho_r^{\mathrm{pw}}&=TKe^{-rT}\1_{\{S_T>K\}}, &\quad
  \rho_y^{\mathrm{pw}}&=-Te^{-rT}\1_{\{S_T>K\}}S_T=-TS\,\Delta^{\mathrm{pw}},\\
  \partial_K^{\mathrm{pw}}&=-e^{-rT}\1_{\{S_T>K\}}, &\quad
  \Theta^{\mathrm{pw}}&=re^{-rT}(S_T-K)^+-e^{-rT}\1_{\{S_T>K\}}S_T\Big(\nu+\frac{\sigma Z}{2\sqrt T}\Big),
\end{aligned}
\end{equation}
where $\rho_r^{\mathrm{pw}}$ combines the discount factor, $-Te^{-rT}(S_T-K)^+$, with the drift,
$Te^{-rT}\1_{\{S_T>K\}}S_T$, and $\Theta^{\mathrm{pw}}$ estimates $\partial_tV=-\partial_TV$. The
sufficient conditions hold for every input when $S,\sigma,T>0$: $f$ is $1$-Lipschitz,
$\Q(S_T=K)=0$ because $S_T$ has a density, and for the input in a small neighbourhood (with
$T$ bounded away from $0$) each derivative of $S_T$ is bounded by an integrable variable of the
form $c_1e^{c_2|Z|}(1+|Z|)$.
So all six are unbiased, as \cref{tab:gr-lab} confirms; the put is the same with
$f'(s)=-\1_{\{s<K\}}$. Two of them are the course's ``pathwise proportionality'':
$\rho_y^{\mathrm{pw}}=-TS\Delta^{\mathrm{pw}}$ and $\rho_r^{\mathrm{pw}}=-TK\partial_K^{\mathrm{pw}}$
hold draw by draw, because the discounted payoff
$\big(Se^{-yT}e^{-\sigma^2T/2+\sigma\sqrt TZ}-Ke^{-rT}\big)^+$ depends on $S$ and $y$ only through
$Se^{-yT}$, and on $K$ and $r$ only through $Ke^{-rT}$ (\exref{ex:gr-homogeneity}). The bumped
$\rho_y$ and $\Delta$ of the course move together for the same reason.

The edges deserve a look. At $S=0$ write $\Delta^{\mathrm{pw}}=e^{-rT}\1_{\{SG>K\}}G$, which is
$0$ for $K>0$, the limit of $e^{-yT}N(d_+)$. For $K=0$ the payoff is the share,
$\Delta^{\mathrm{pw}}=e^{-rT}G$ has mean $e^{-yT}$ and variance $e^{-2yT}(e^{\sigma^2T}-1)$, small
when $\sigma^2T$ is. As $\sigma\sqrt T\to0$ the hypothesis $\Q(S_T=K)=0$ survives for every
$\sigma>0$, but at $\sigma=0$ with $K$ equal to the forward it fails, and indeed the price
$(Se^{-yT}-Ke^{-rT})^+$ has a kink there. As $T\downarrow0$ the theta estimator contains
$\sigma Z/(2\sqrt T)$. At and in the money its variance grows like $1/T$ ($T\Var\Theta^{\mathrm{pw}}$
tends to $\frac14S^2\sigma^2\Var(Z^+)=34.1$ at $S=100$ and to $\frac14S^2\sigma^2=121$ at $S=110$), while
out of the money it tends to $0$ ($37.6$, $0.019$ at $T=1$, $0.01$ for $S=90$); the theta itself grows
like $1/\sqrt T$ at the money and stays bounded away from it.

\begin{remark}[Why digitals fail, and three remedies]\label{rem:gr-digital}
For the digital, $Y(S)=e^{-rT}\1_{\{SG>K\}}$ is differentiable at $S$ with derivative $0$ for
almost every draw, so the pathwise estimator is identically zero, while the delta is
$0.0195$. Condition (PW2) is what fails. For a draw with critical spot $K/G$ at distance
$u=|S-K/G|<\delta$ from $S$, the quotient $(Y(S')-Y(S))/(S'-S)$ equals $e^{-rT}/|S'-S|$ for
$S'$ just across $K/G$, so any admissible $\kappa$ must exceed $e^{-rT}/u$; and $u$ has a
density bounded below near $0$, so $\E[1/u;u<\delta]=\infty$. Indeed no integrable $\kappa$ can
exist, since otherwise the theorem would give the delta $0$. The same happens for barrier
options, whose payoff jumps when the path touches the barrier, and for the pathwise derivative
of the call's pathwise delta, which jumps at the strike (\cref{sec:gr-gamma}). Three remedies:
\begin{enumerate}[label=(\alph*)]
  \item \emph{Smoothing}: replace the payoff by a Lipschitz approximation, such as the call
    spread $\big((S_T-K+\eta)^+-(S_T-K)^+\big)/\eta$, and accept a bias of order $\eta$;
    the variance of its pathwise delta grows like $1/\eta$, the trade-off of
    \cref{prop:gr-finite-difference}(iv) again.
  \item \emph{Likelihood ratios}, which differentiate the density instead of the payoff
    (\cref{sec:gr-lr}).
  \item \emph{Conditional expectation}: replace $Y$ by $\E[Y\mid\mathcal G]$ for a
    $\sigma$-algebra $\mathcal G$ such that the conditional payoff is a smooth function of the
    input; this keeps the mean, lowers the variance (\cref{thm:st-conditional}) and makes the
    pathwise method valid. For a payoff on a simulated path, conditioning on everything but the
    last step replaces the indicator of the last step by a Black--Scholes digital price
    (\exref{ex:gr-smoothing}).
\end{enumerate}
\end{remark}

\paragraph{Along a path.} For an Asian call on the arithmetic average
$A=\frac1m\sum_jS_{t_j}$ of $m$ equally spaced fixings, $t_j=jT/m$, simulated exactly with
\texttt{paths.gbm\_paths} (\cref{prop:ps-gbm-exact}) from $S_{t_j}=S\exp(\nu t_j+\sigma W_{t_j})$,
\begin{equation}\label{eq:gr-asian-pw}
  \Delta^{\mathrm{pw}}=e^{-rT}\1_{\{A>K\}}\frac AS,\qquad
  \mathcal V^{\mathrm{pw}}=e^{-rT}\1_{\{A>K\}}\frac1m\sum_{j=1}^mS_{t_j}\big(W_{t_j}-\sigma t_j\big).
\end{equation}
The sufficient conditions hold as for the call ($A$ has a density). For the geometric average
$A^{\mathrm g}$ the same formulas, with $\partial_\sigma A^{\mathrm g}=A^{\mathrm g}\cdot\frac1m\sum_j(W_{t_j}-\sigma t_j)$, can be checked
against the closed form (\cref{prop:as-geometric}): with $m=12$ the exact delta and vega are
$0.486882$ and $22.4010$, and $4\times10^5$ paths give $0.4866(9)$ and $22.42(6)$ (our
implementation, seed $3$; standard errors in units of the last digit). The arithmetic Asian
gives $0.4951(9)$ and $24.07(7)$ from the same paths.

When the path is produced by a discretisation scheme, the pathwise derivative is propagated
along with it. For the Euler scheme $X_{k+1}=X_k+a(X_k)h+b(X_k)\Delta W_k$ of
\cref{thm:ps-euler}, the derivative $D_k=\partial X_k/\partial x_0$ obeys
\begin{equation}\label{eq:gr-tangent}
  D_{k+1}=D_k\big(1+a'(X_k)h+b'(X_k)\Delta W_k\big),\qquad D_0=1,
\end{equation}
and $\partial_{x_0}f(X_N)=f'(X_N)D_N$. If $|a'|,|b'|\le L$ and $f$ is Lipschitz, each step is
Lipschitz in its input with constant $1+Lh+L|\Delta W_k|$, so (PW2) holds with
$\kappa=L_f\prod_k(1+Lh+L|\Delta W_k|)$, whose mean is the product of the means of independent
factors. If $b$ does not vanish, $X_N$ has a density (given $X_{N-1}$ it is Gaussian), so (PW1)
holds for every $f$ that is differentiable except at finitely many points, such as calls, puts
and their combinations. For geometric Brownian motion,
$D_N=X_N/x_0$ exactly. This is \emph{forward mode}: one derivative per input, computed
alongside the path. \Cref{sec:gr-aad} runs it backwards.

\section{The likelihood ratio method}\label{sec:gr-lr}

The pathwise method moves the payoff and keeps the law. The likelihood-ratio method keeps the
payoff and moves the law. Let $X$ have density $g_\theta$ on $\R^k$, let the payoff $f$ not
depend on $\theta$, and write $\alpha(\theta)=\E_\theta[f(X)]=\int f\,g_\theta\,\dd x$.

\begin{theorem}[Likelihood ratio method]\label{thm:gr-likelihood-ratio}
Let $\theta_0\in I$, an open interval, $f:\R^k\to\R$ Borel, and assume:
\begin{enumerate}[label=(LR\arabic*)]
  \item there is $\delta>0$ such that, for every $x$, $\theta\mapsto g_\theta(x)$ is
    differentiable on $U=(\theta_0-\delta,\theta_0+\delta)$, with derivative
    $\dot g_\theta(x)$;
  \item there is a Borel $H\ge0$ with $|\dot g_\theta(x)|\le H(x)$ for all $\theta\in U$ and $x$, and
    $\int_{\R^k}|f(x)|H(x)\,\dd x<\infty$.
\end{enumerate}
If $\E_{\theta_0}|f(X)|<\infty$, then $\alpha$ is differentiable at $\theta_0$ and
\[
  \alpha'(\theta_0)=\E_{\theta_0}\big[f(X)\,\ell_{\theta_0}(X)\big],\qquad
  \ell_\theta(x)=\frac{\dot g_\theta(x)}{g_\theta(x)}\ \text{ on }\{g_\theta>0\},\quad\ell_\theta=0\ \text{elsewhere}.
\]
The \emph{score} $\ell_{\theta_0}$ does not involve $f$, so one simulation under $g_{\theta_0}$
gives every Greek of every payoff; the estimator $f(X)\ell_{\theta_0}(X)$ has finite variance
if and only if $\E_{\theta_0}[f(X)^2\ell_{\theta_0}(X)^2]<\infty$.
\end{theorem}

\begin{proof}
For $0<|h|<\delta$ and every $x$, the mean value theorem gives
$|g_{\theta_0+h}(x)-g_{\theta_0}(x)|/|h|\le H(x)$, so the integrands
$f(x)\big(g_{\theta_0+h}(x)-g_{\theta_0}(x)\big)/h$ are dominated by the integrable $|f|H$ and
converge to $f(x)\dot g_{\theta_0}(x)$. By dominated convergence,
$\alpha'(\theta_0)=\int f\dot g_{\theta_0}\,\dd x$. Where $g_{\theta_0}(x)=0$, the nonnegative
function $\theta\mapsto g_\theta(x)$ has a minimum at $\theta_0$, so $\dot g_{\theta_0}(x)=0$;
elsewhere $\dot g_{\theta_0}=\ell_{\theta_0}g_{\theta_0}$. Hence
$\int f\dot g_{\theta_0}\,\dd x=\int f\ell_{\theta_0}g_{\theta_0}\,\dd x$.
\end{proof}

Condition (LR2) is a domination condition. It makes the difference quotients
\[
  f(X)\Big(\frac{g_{\theta_0+h}(X)}{g_{\theta_0}(X)}-1\Big)\Big/h,\qquad 0<|h|<\delta,
\]
uniformly integrable under $g_{\theta_0}$. When the support $\{g_\theta>0\}$ does not depend on
$\theta$, these quotients have mean $(\alpha(\theta_0+h)-\alpha(\theta_0))/h$, and uniform
integrability together with almost sure convergence is all that the interchange needs
\citep[Section~4.6]{durrett}. A moving support is not in itself an obstacle: a location family
$\psi(x-\theta)$ with $\psi$ a smooth, compactly supported density satisfies (LR1)--(LR2). The
uniform law on $[0,\theta]$, on the other hand, violates (LR1): its density is not differentiable
in $\theta$ at $x=\theta$.

The likelihood ratio method is importance sampling read the other way round
(\cref{ch:09-importance-sampling}). If $g_{\theta_0}>0$ wherever $g_\theta>0$, then
$\alpha(\theta)=\E_{\theta_0}\big[f(X)g_\theta(X)/g_{\theta_0}(X)\big]$ by
\cref{prop:is-likelihood-ratio}, with sampling law $g_{\theta_0}$ and target $g_\theta$; the
score is the derivative of that weight at $\theta=\theta_0$.

\paragraph{Black--Scholes.} The log-price $\log S_T$ has the density of
$\mathcal N(\log S+\nu T,\sigma^2T)$. Differentiating its logarithm in each input and writing
the result in terms of $Z=(\log S_T-\log S-\nu T)/(\sigma\sqrt T)$, the estimators are, with
$Y=e^{-rT}f(S_T)$ for any payoff $f$,
\begin{equation}\label{eq:gr-lr-bs}
\begin{aligned}
  \Delta^{\mathrm{lr}}&=Y\frac{Z}{S\sigma\sqrt T}, &\quad
  \Gamma^{\mathrm{lr}}&=Y\frac{Z^2-1-\sigma\sqrt TZ}{S^2\sigma^2T}, &\quad
  \mathcal V^{\mathrm{lr}}&=Y\Big(\frac{Z^2-1}{\sigma}-\sqrt TZ\Big),\\
  \rho_y^{\mathrm{lr}}&=-Y\frac{\sqrt TZ}{\sigma}=-TS\Delta^{\mathrm{lr}}, &\quad
  \rho_r^{\mathrm{lr}}&=-TY+Y\frac{\sqrt TZ}{\sigma}, &\quad
  \Theta^{\mathrm{lr}}&=rY-Y\Big(\frac{\nu Z}{\sigma\sqrt T}+\frac{Z^2-1}{2T}\Big).
\end{aligned}
\end{equation}
The gamma uses $\partial_S^2g/g$, computed in \cref{sec:gr-gamma}; $\rho_r^{\mathrm{lr}}$ adds the
discount factor's $-TY$ to the score of the drift. For the digital,
$\Delta^{\mathrm{lr}}=e^{-rT}\1_{\{S_T>K\}}Z/(S\sigma\sqrt T)$ is unbiased with per-sample
standard deviation $0.0290$ (\cref{exa:gr-digital}). The domination (LR2) holds for every
payoff of polynomial growth in $s$ and $1/s$: for the mean, for instance,
$\sup_{|m'-m|<\delta}|\partial_mq(x-m')|\le C(1+|x-m|)e^{-((|x-m|-\delta)^+)^2/(2\sigma^2T)}$ for
the Gaussian density $q$.

\begin{table}[ht]
  \centering
  \small
  \begin{tabular}{@{}lccccrr@{}}
    \toprule
    & Black--Scholes & pathwise & likelihood ratio & bump, CRN & \multicolumn{2}{c}{sd per sample}\\
    & & & & & pathwise & LR \\
    \midrule
    $\Delta$ & $0.509643$ & $0.5108(18)$ & $0.5093(41)$ & $0.5108(18)$ & $0.582$ & $1.287$ \\
    $\Gamma$ & $0.0195277$ & $0.01954(11)^\ast$ & $0.01941(39)$ & $0.01977(36)$ & $0.0363^\ast$ & $0.1234$ \\
    $\mathcal V$ & $39.0554$ & $39.08(23)$ & $38.82(79)$ & $39.08(23)$ & $72.67$ & $246.70$ \\
    $\rho_r$ & $43.6000$ & $43.71(16)$ & $43.55(37)$ & $43.71(16)$ & $49.15$ & $117.08$ \\
    $\rho_y$ & $-50.9643$ & $-51.08(18)$ & $-50.93(41)$ & $-51.08(18)$ & $58.22$ & $128.71$ \\
    $\partial_tV$ & $-3.32226$ & $-3.324(21)$ & $-3.299(75)$ & $-3.324(21)$ & $6.73$ & $23.45$ \\
    \bottomrule
  \end{tabular}
  \caption{The six Greeks of the test-case call from the same $n=10^5$ draws (our
    implementation, \texttt{default\_rng}, seed $13$, \texttt{greeks.greeks\_table}); standard
    errors in units of the last digit. Bumps with common random numbers: $1\%$ of $S$ (second
    difference for $\Gamma$), of $\sigma$ and of $T$, and one basis point for $r$ and $y$.
    The per-sample standard deviations are exact (quadrature). $^\ast$For the gamma, the
    PW-LR estimator of \cref{prop:gr-hybrid-gamma}; its LR-PW estimator gives $0.019557(92)$,
    with standard deviation $0.0290$.}
  \label{tab:gr-lab}
\end{table}

\Cref{tab:gr-lab} compares the three methods on the test case. Every estimate is within two
standard errors of the closed form. The bumps with common random numbers reproduce the pathwise
estimates and their errors, as \cref{thm:gr-pathwise} predicts. The likelihood-ratio
estimators have standard deviations $2.2$ to $3.5$ times larger, that is $5$ to $12$ times more
paths for the same accuracy; the course's observation that the likelihood-ratio $\rho_y$ has
errors ``more than twice'' the pathwise ones is the ratio $128.71/58.22=2.21$; the exact ratio
grows with the spot, from $1.86$ at $S_0=90$ to $2.81$ at $S_0=110$, and exceeds $2$ for
$S_0\gtrsim95$. In
\cref{fig:gr-course}(c) the pathwise delta has the $1/n$ steps of \cref{rem:gr-kernel}, while
the likelihood-ratio delta is smooth in $S_0$, since each draw contributes the smooth
$e^{-rT}(SG-K)^+Z/(S\sigma\sqrt T)$, but further from the truth.

\begin{remark}[The weight degenerates as $\sigma\sqrt T\to0$]\label{rem:gr-lr-degenerate}
The delta weight $Z/(S\sigma\sqrt T)$ has standard deviation $1/(S\sigma\sqrt T)$. Let $f$ be
continuous at the forward $F=Se^{(r-y)T}$, with $|f(s)|\le C(1+s^p)$. Then, as $\sigma\downarrow0$
with $S,T,r,y$ fixed,
\[
  \sigma^2T\,\Var\big(\Delta^{\mathrm{lr}}\big)\longrightarrow e^{-2rT}f(F)^2/S^2 .
\]
Indeed $\sigma^2T\Var(\Delta^{\mathrm{lr}})=e^{-2rT}\E[f(S_T)^2Z^2]/S^2-\big(\sigma\sqrt T\,\Delta\big)^2$;
as $\sigma\downarrow0$, $S_T\to F$ for every $Z$, and for $\sigma\le1$ the integrand is
dominated by $C'(1+e^{2p\sqrt T|Z|})Z^2$, with $C'$ depending on $S,T,r,y,p$, so the first term converges to
$e^{-2rT}f(F)^2/S^2$ by dominated convergence; the same argument gives
$\sigma\sqrt T\Delta=e^{-rT}\E[f(S_T)Z]/S\to e^{-rT}f(F)\E[Z]/S=0$. So whenever $f(F)\ne0$ the
standard deviation grows like $1/(\sigma\sqrt T)$, while the pathwise delta of a payoff with
Lipschitz constant $L_f$ has variance at most $e^{-2rT}L_f^2\E[G^2]$, bounded. For the share
($K=0$) the two exact variances are $e^{-2yT}\big(e^{\sigma^2T}(1+4\sigma^2T)/(\sigma^2T)-1\big)=28.04$
and $e^{-2yT}(e^{\sigma^2T}-1)=0.0392$ at the test case, a ratio of $715$. The same happens as
$T\downarrow0$, and the gamma weight is worse, of order $1/(\sigma^2T)$. At the forward,
$f(F)=0$ for the call and the limit is finite (\exref{ex:gr-small-vol}).
\end{remark}

\paragraph{Along a path.} For a payoff of the whole path, the law of the path is a product
of Gaussian steps: with $\Delta t=T/m$ and $\log(S_{t_j}/S_{t_{j-1}})=\nu\Delta t+\sigma\sqrt{\Delta t}Z_j$,
the input $S$ enters only the first step and $\sigma$ enters every step, so the scores are
\begin{equation}\label{eq:gr-asian-lr}
  \Delta^{\mathrm{lr}}=Y\frac{Z_1}{S\sigma\sqrt{t_1}},\qquad
  \mathcal V^{\mathrm{lr}}=Y\sum_{j=1}^m\Big(\frac{Z_j^2-1}{\sigma}-\sqrt{\Delta t}\,Z_j\Big),
\end{equation}
valid for any payoff of the fixings, digital or barrier included. Both degrade as the grid is
refined (\cref{fig:gr-asian}(a)). The vega weight is a sum of $m$ independent terms of variance
$2/\sigma^2+\Delta t$, so its variance grows linearly in $m$; in continuous time the laws of the
paths for two volatilities are mutually singular, because the quadratic variation of
$\log S$ reveals $\sigma$, and there is no likelihood-ratio vega at all. The delta weight has
standard deviation $1/(S\sigma\sqrt{t_1})$, which grows as the first fixing approaches $0$. This
weight is Girsanov's (\cref{thm:is-girsanov}): moving $\log S$ by $\epsilon$ moves every fixing,
and the cheapest drift that achieves this by time $t_1$ is the constant $\epsilon/(\sigma t_1)$ on
$[0,t_1]$ (by Cauchy--Schwarz it has the smallest $\int\theta^2$ among drifts with
$\int_0^{t_1}\theta=\epsilon/\sigma$, hence the smallest second moment of the weight),
and its weight differentiates to $W_{t_1}/(S\sigma t_1)=Z_1/(S\sigma\sqrt{t_1})$. For
$m=1$ the Asian is the call, and the likelihood-ratio standard deviations are $2.2$ (delta) and
$3.4$ (vega) times the pathwise ones (\cref{tab:gr-lab}); for the arithmetic Asian with $m=128$
they are $8.9$ and $17$ times ($10^5$ paths, seed $2026$).

\begin{figure}[ht]
  \centering
  \includegraphics[width=0.94\textwidth]{ch13-asian}
  \caption{Asian call on the test case. (a) Per-sample standard deviations of the pathwise
    \eqref{eq:gr-asian-pw} and likelihood-ratio \eqref{eq:gr-asian-lr} deltas and vegas
    against the number of fixings, divided by the pathwise value at $m=1$ ($10^5$ paths per $m$,
    seed $2026$). (b) Bucketed vegas $\partial P/\partial\sigma_j$ for $m=12$ from one adjoint
    sweep (bars, with $95\%$ intervals) and from $24$ central bumps of $10^{-4}$ with common
    random numbers (dots), same $10^5$ paths (seed $2027$); the total is the parallel vega.}
  \label{fig:gr-asian}
\end{figure}

\section{The gamma}\label{sec:gr-gamma}

A second derivative is harder. For the call, the pathwise delta
$e^{-rT}\1_{\{SG>K\}}G$ jumps when the critical spot $K/G$ crosses $S$, so its own pathwise
derivative is $0$ almost everywhere: the ``pathwise gamma'' is identically zero
(\cref{rem:gr-digital}), and the second difference with common draws is a kernel estimator with
variance of order $1/(nh)$ (\cref{rem:gr-kernel}). The likelihood-ratio gamma of
\eqref{eq:gr-lr-bs} is unbiased but noisy. The course's way out is to differentiate twice by
different methods: once pathwise, where the payoff is smooth enough, and once by likelihood
ratio, where it is not.

\begin{proposition}[Mixed estimators of the gamma]\label{prop:gr-hybrid-gamma}
Let $S,\sigma,T>0$, $S_T=SG$ as above, and let $f:(0,\infty)\to\R$ be continuous, and
continuously differentiable except at finitely many points, with $|f(s)|+|sf'(s)|\le C(s^p+s^{-p})$
wherever $f'(s)$ exists, for some constants $C,p$.
Then $V(S)=e^{-rT}\E f(S_T)$ is twice differentiable, and its gamma $\Gamma=V''(S)$ is the
expectation of each of
\[
\begin{aligned}
  &\text{(LR-LR)}\quad e^{-rT}f(S_T)\,\frac{Z^2-1-\sigma\sqrt TZ}{S^2\sigma^2T},\\
  &\text{(PW-LR)}\quad e^{-rT}f'(S_T)S_T\,\frac{Z-\sigma\sqrt T}{S^2\sigma\sqrt T},\\
  &\text{(LR-PW)}\quad e^{-rT}\big(f'(S_T)S_T-f(S_T)\big)\frac{Z}{S^2\sigma\sqrt T},
\end{aligned}
\]
all with finite variance. PW-LR applies the likelihood ratio to the pathwise delta, LR-PW
differentiates the likelihood-ratio delta pathwise. The LR-LR estimator needs only $f$ Borel
with $|f(s)|\le C(s^p+s^{-p})$, digitals included; the other two are biased for the digital.
\end{proposition}

\begin{proof}
Let $u=\log S$, $X=\nu T+\sigma\sqrt TZ$, so that $S_T=e^{u+X}$, and $q$ the density of $X$.
For a Borel $k$ with $|k(s)|\le C(s^p+s^{-p})$ let $\phi_k(u)=\E[k(e^{u+X})]=\int k(e^x)q(x-u)\,\dd x$.
The family $u\mapsto q(x-u)$ satisfies (LR1) and (LR2), and so does $u\mapsto\partial_uq(x-u)$ in
place of $q(x-u)$, since $\partial_u^jq(x-u)$ is a polynomial in $x-u$ times $q(x-u)$ and $k(e^x)$
grows at most like $e^{p|x|}$. By \cref{thm:gr-likelihood-ratio}, and the same dominated
convergence argument for the second derivative,
\[
  \phi_k'(u)=\frac{\E[k(S_T)Z]}{\sigma\sqrt T},\qquad \phi_k''(u)=\frac{\E[k(S_T)(Z^2-1)]}{\sigma^2T},
\]
because $\partial_u\log q(x-u)=(x-u-\nu T)/(\sigma^2T)=Z/(\sigma\sqrt T)$ and
$\partial_u^2q/q=(Z^2-1)/(\sigma^2T)$. Next, $f$ is continuous and piecewise $C^1$, so
$f(e^{u'+X})-f(e^{u+X})=\int_u^{u'}f'(e^{v+X})e^{v+X}\,\dd v$, and
$|f(e^{u'+X})-f(e^{u+X})|\le|u'-u|\,\kappa$ for $|u'-u|<1$, with
$\kappa=\sup_{|v-u|<1}C(e^{p(v+X)}+e^{-p(v+X)})$ integrable, also after multiplication by $|Z|$;
and $f$ is differentiable at $S_T$ almost surely, since $S_T$ has a density and $f$ fails to be
differentiable at finitely many points only. By \cref{thm:gr-pathwise},
$\phi_f'(u)=\E[f'(S_T)S_T]$ and $\frac{\dd}{\dd u}\E[f(S_T)Z]=\E[f'(S_T)S_TZ]$.

Since $\partial_S=S^{-1}\partial_u$, $\Gamma=e^{-rT}S^{-2}(\phi_f''-\phi_f')(u)$. Writing
$\phi_f''-\phi_f'$ with both derivatives by likelihood ratio gives LR-LR. Writing
$\phi_f'=\phi_k$ with $k(s)=sf'(s)$ (pathwise) and then
$\phi_f''=\phi_k'=\E[k(S_T)Z]/(\sigma\sqrt T)$ (likelihood ratio, which needs only the growth
of $k$) gives PW-LR. Writing $\phi_f'=\E[f(S_T)Z]/(\sigma\sqrt T)$ (likelihood ratio) and
differentiating it pathwise gives $\phi_f''=\E[f'(S_T)S_TZ]/(\sigma\sqrt T)$, hence LR-PW. Each
estimator is a polynomial in $Z$ times a function of $S_T$ of polynomial growth in $S_T$ and
$1/S_T$, so it has finite variance. For the digital, $f$ is not continuous and $f'=0$ off $K$:
PW-LR is then identically $0$ and LR-PW has mean $-\Delta/S$.
\end{proof}

For the test-case call the per-sample standard deviations are $0.1234$ (LR-LR), $0.0363$
(PW-LR) and $0.0290$ (LR-PW): the mixed estimators need $12$ and $18$ times fewer paths than
the pure likelihood ratio. The second difference with common draws and a $1\%$ bump has standard
deviation $0.112$ and a bias of $0.02\%$: comparable to LR-LR, at three evaluations per draw. For the call,
$f'(s)s-f(s)=K\1_{\{s>K\}}$, and LR-PW is $K/S$ times the likelihood-ratio delta of the digital:
the gamma of a call is the delta of a digital in disguise. \Cref{fig:gr-course}(d) is the
course's picture: with $1024$ draws the mixed gamma is already smooth in $S_0$, and LR-LR
converges as well, more slowly. Both mixed estimators inherit the likelihood-ratio weight, so
they degenerate like $1/(\sigma\sqrt T)$, not like the $1/(\sigma^2T)$ of LR-LR.

\demo{13-mc-greeks.ipynb}{pathwise, likelihood-ratio and mixed Greeks against the spot, for
any number of draws.}

\section{Many Greeks at once: adjoints}\label{sec:gr-aad}

A book may depend on hundreds of inputs: a volatility per maturity, a forward rate per period.
Bumping costs one or two repricings per input, and so does the pathwise method in forward mode
\eqref{eq:gr-tangent}, one derivative at a time. The \emph{adjoint}, or reverse, mode computes
the derivative of one output with respect to all the inputs in one backward sweep.

\begin{remark}[Adjoint algorithmic differentiation]\label{rem:gr-aad}
Write the computation of a path's discounted payoff $P$ from the inputs $x_1,\dots,x_p$ as a
sequence of elementary operations $v_k=\varphi_k(v_j:j\prec k)$, each using earlier values.
Forward mode propagates tangents $\dot v_k=\sum_{j\prec k}\partial_j\varphi_k\,\dot v_j$ from a
seed $\dot x$, one input direction per sweep. Reverse mode propagates adjoints
$\bar v_j=\partial P/\partial v_j$ backwards, $\bar v_j=\sum_{k\succ j}\bar v_k\,\partial_j\varphi_k$,
from $\bar P=1$; at the end, $\bar x_i=\partial P/\partial x_i$ for every $i$. Each elementary
operation is visited once in each sweep, and its local partial derivatives cost a bounded
multiple of the operation itself. In the operation-count model of
\citet[Section~4.6]{griewank-walther}, the price and its whole gradient together cost at most
four times the price alone, whatever $p$. \textcite{giles-glasserman-adjoints} quote this bound and use the adjoint to compute the
deltas and vegas of the LIBOR market model, with an advantage over forward mode that grows in
proportion to the number of forward rates; \textcite{capriotti-aad} develops the method for
pricing code in general. Three caveats. (a) The backward sweep needs the intermediate values of the forward
sweep, stored (the ``tape'') or recomputed; in Monte Carlo this is one path at a time, so the
memory is proportional to the number of time steps. (b) Applied path by path, the adjoint
computes \emph{exactly} the pathwise estimator, to rounding: it inherits the hypotheses of
\cref{thm:gr-pathwise}, so digitals and barriers must be smoothed first, and a gamma needs a
second-order method such as those of \cref{prop:gr-hybrid-gamma}. (c) The bound is on
operations; actual running times also depend on memory traffic and on the implementation.
\end{remark}

For the Asian call with a volatility $\sigma_j$ on each period $(t_{j-1},t_j]$, the forward sweep
is $X_0=\log S$, $X_j=X_{j-1}+(r-y-\sigma_j^2/2)\Delta t+\sigma_j\sqrt{\Delta t}Z_j$,
$S_{t_j}=e^{X_j}$, $A=\frac1m\sum_jS_{t_j}$, $P=e^{-rT}(A-K)^+$. Differentiating each line
backwards gives $\bar A=e^{-rT}\1_{\{A>K\}}$ and, for $j=m,\dots,1$,
\[
  \bar X_j=\bar A\,\frac{S_{t_j}}m+\bar X_{j+1}\quad(\bar X_{m+1}=0),\qquad
  \frac{\partial P}{\partial\sigma_j}=\bar X_j\big(\sqrt{\Delta t}Z_j-\sigma_j\Delta t\big),
\]
then $\Delta=\bar X_1/S$, $\rho_y=-\Delta t\sum_j\bar X_j$ and
$\rho_r=-TP+\Delta t\sum_j\bar X_j$: all $m+3$ sensitivities for one forward and one backward
pass, against $2(m+3)$ repricings for central bumps. With a flat volatility the bucketed vegas
add up to the pathwise vega \eqref{eq:gr-asian-pw}. In our implementation the adjoint agrees
with the $m+3$ forward-mode sweeps to $10^{-12}$ on every path and with central bumps of
$10^{-4}$ to $10^{-3}$, and \cref{fig:gr-asian}(b) shows the twelve bucketed vegas: the early periods
matter most, because a volatility early in the year moves every later fixing.

\begin{incode}
\begin{lstlisting}
x = np.full(n, np.log(S))                       # forward sweep, storing the tape X
for j in range(m):
    x = x + (r - y - sig[j]**2 / 2) * dt + sig[j] * sq * Z[:, j]
    X[:, j] = x
Sj = np.exp(X); A = Sj.mean(axis=1); P = disc * np.maximum(A - K, 0.0)
Abar, acc = disc * (A > K), 0.0                 # reverse sweep
for j in reversed(range(m)):
    acc = acc + Abar * Sj[:, j] / m
    Xbar[:, j] = acc
vega = Xbar * (sq * Z - sig * dt); delta = Xbar[:, 0] / S
\end{lstlisting}
\texttt{asian\_adjoint} in \coderef{code/src/mfe230d/greeks.py}; \texttt{asian\_tangent} is the
forward mode, one input direction per sweep. No automatic-differentiation library is used:
the backward sweep is the chain rule written by hand.
\end{incode}

\demo{13-mc-greeks.ipynb}{Asian Greeks against the number of fixings, and bucketed vegas by
adjoint and by bumping, with timings.}

\furtherreading
\Cref{fig:gr-course} and \cref{tab:gr-lab} reproduce the course's bumping study and its
analytical-differencing, likelihood-ratio and mixed estimators. \textcite{glasserman-mc},
Chapter~7, is the reference for everything here: finite differences in Section~7.1, the
pathwise method in Section~7.2, likelihood ratios and the gamma in Section~7.3.
\textcite{broadie-glasserman-greeks} introduced the pathwise and likelihood-ratio estimators
to option pricing, and \citet[Chapter~VII]{asmussen-glynn} treat derivative estimation in
simulation in general. The Malliavin-calculus weights of \textcite{fournie-malliavin} generalise
the likelihood ratio to models without explicit densities. Adjoints are due to
\textcite{giles-glasserman-adjoints} in this context; \textcite{capriotti-aad} and
\textcite{griewank-walther} cover algorithmic differentiation in depth.
